% 01_problema.tex — Definición del problema 3DM
\section{Definición del problema}
\label{sec:problema}

\subsection{Descripción informal}
\label{sec:problema:informal}

El problema \threedm{} (\textbf{3DM}) generaliza el matching bipartito clásico
a tres dimensiones.  Dado que el matching bipartito ---que empareja pares
$(x,y)\in X\times Y$--- se resuelve en tiempo polinomial mediante el algoritmo
de Hopcroft--Karp~\cite{hopcroft1973}, es natural preguntarse qu\'e ocurre al
agregar una \emph{tercera coordenada}.  La respuesta es categ\'orica: el
problema se vuelve \textbf{NP-completo} (Karp~\cite{karp1972}, problema~\#17).

\subsection{Definición formal}
\label{sec:problema:formal}

\begin{definition}[Instancia de 3DM]
Una instancia de \threedm{} es una 4-tupla
$I = (X,\,Y,\,Z,\,T)$
donde $X$, $Y$, $Z$ son conjuntos finitos \emph{disjuntos} con
$|X|=|Y|=|Z|=n$, y
$T \subseteq X\times Y\times Z$ es el conjunto de \emph{tripletas candidatas},
con $m = |T|$.
\end{definition}

\begin{definition}[3-matching]
Un subconjunto $M\subseteq T$ es un \emph{3-matching} si para todo par de
tripletas distintas $(x,y,z),(x',y',z')\in M$ se cumple
$x\neq x'$, $y\neq y'$ y $z\neq z'$
(cada elemento de $X\cup Y\cup Z$ aparece a lo más una vez en $M$).
\end{definition}

\paragraph{Versión de decisión (NP-completa).}
¿Existe un 3-matching \emph{perfecto} de tamaño $n$?
Equivalentemente: ¿existen $n$ tripletas en $T$ cuyas proyecciones sobre $X$,
$Y$ y $Z$ cubren exactamente $X$, $Y$ y $Z$ sin repetición?

\paragraph{Versión de optimización (max-3DM, NP-difícil).}
\[
  \OPT(I) \;=\; \max_{M\subseteq T,\;M\text{ 3-matching}} |M|.
\]
La versión de decisión equivale a comprobar si $\OPT(I)=n$.
Este trabajo implementa la versión de \textbf{optimización}, ya que un solver
de optimización resuelve gratis la decisión y produce información más rica para
el análisis experimental.

\subsection{Ejemplo ilustrativo}
\label{sec:problema:ejemplo}

Tomemos $n=3$, $X=\{x_1,x_2,x_3\}$, $Y=\{y_1,y_2,y_3\}$,
$Z=\{z_1,z_2,z_3\}$, y el conjunto de $m=5$ tripletas:
\[
T = \bigl\{
  t_1{=}(x_1,y_1,z_1),\;
  t_2{=}(x_1,y_2,z_3),\;
  t_3{=}(x_2,y_2,z_2),\;
  t_4{=}(x_3,y_3,z_3),\;
  t_5{=}(x_3,y_1,z_2)
\bigr\}.
\]
El matching óptimo es
$M^* = \{t_1, t_3, t_4\}$, con $|M^*|=3=n$ (matching perfecto).
El candidato $M' = \{t_2, t_3, t_5\}$ tiene también tamaño~3 pero es inválido:
$y_2$ aparece en $t_2$ y en $t_3$, y $z_2$ aparece en $t_3$ y en $t_5$.

\subsection{Aplicaciones}
\label{sec:problema:aplicaciones}

\threedm{} tiene aplicaciones naturales en escenarios donde tres recursos
distintos deben asignarse simultáneamente:
\begin{itemize}
  \item \textbf{Asignación tripartita académica}: emparejar estudiantes,
        proyectos y tutores respetando compatibilidades tripartitas.
  \item \textbf{Scheduling con tres recursos no-fungibles}: turnos de
        quirófano (sala~+~equipo médico~+~paciente), calendarios deportivos
        (equipo~+~cancha~+~árbitro).
  \item \textbf{Bioinformática}: alineación tripartita de lecturas genómicas,
        posiciones y muestras experimentales.
  \item \textbf{Logística \textit{last-mile}}: asignación paquete--vehículo--conductor
        con restricciones de peso, licencia y ruta.
\end{itemize}

\subsection{Versiones del problema}
\label{sec:problema:versiones}

\begin{table}[htbp]
\centering
\caption{Versiones de 3DM y su complejidad.}
\label{tab:versiones}
\begin{tabular}{lp{5.5cm}l}
\toprule
Versión & Pregunta & Complejidad \\
\midrule
Decisión (3DM) & ¿Existe $M\subseteq T$ con $|M|=n$? & NP-completo \cite{karp1972} \\
Optimización (max-3DM) & Maximizar $|M|$ & NP-difícil; APX-difícil \cite{kann1991} \\
Conteo (\#3DM) & ¿Cuántos matchings perfectos existen? & \#P-completo \cite{valiant1979} \\
\bottomrule
\end{tabular}
\end{table}
